\chapter{Methodology}
\label{chap:methodology}

\section{Overview}
\label{sec:methodology_overview}

This chapter describes the design, implementation, and evaluation methodology of
Block-Insure, a blockchain-backed prototype for health-insurance claim
management. The prototype combines an on-chain claim lifecycle with off-chain
evidence storage, synthetic healthcare-registry verification, a two-oracle
quorum, explainable fraud-risk scoring, auditor voting, and on-chain settlement.

The research follows a design-and-evaluation methodology. First, a working
artifact was designed to address four recurring insurance-claim problems:
unauditable state changes, duplicate evidence, dependence on a single external
verifier, and opaque fraud decisions. Second, the artifact was implemented as a
full-stack prototype. Third, its principal technical claims were evaluated using
automated contract tests, code coverage, a gas-cost comparison, a held-out
fraud-model evaluation, and an exploratory concurrency benchmark.

The methodology deliberately distinguishes between \emph{enforced controls} and
\emph{audit signals}. Enforced controls, such as role checks, duplicate
detection, state-transition restrictions, oracle quorum, and settlement
preconditions, are executed by the smart contract. Audit signals, such as the
off-chain evidence chain, registry inclusion proof, oracle response hash, and
reserve-intelligence dashboard, improve traceability but do not independently
prevent a privileged operator from supplying incorrect off-chain data.

\begin{table}[htbp]
\centering
\caption{Research questions and corresponding evaluation methods}
\label{tab:research_questions}
\resizebox{\textwidth}{!}{
\begin{tabular}{|p{1.2cm}|p{7.0cm}|p{6.0cm}|}
\hline
\textbf{RQ} & \textbf{Research Question} & \textbf{Evaluation Method} \\
\hline
RQ1 & Can the prototype enforce an auditable insurance-claim lifecycle and
prevent unauthorized or invalid state transitions? &
Hardhat functional tests and Solidity code coverage \\
\hline
RQ2 & Does committing one registry Merkle root reduce on-chain gas consumption
compared with storing every registry-record hash? &
Controlled gas comparison at six registry sizes \\
\hline
RQ3 & Can the explainable Bayesian scorer separate the deliberately constructed
fraud scenarios from legitimate records in the synthetic registry? &
Deterministic training/held-out split, baselines, confusion matrix, ROC AUC, and
Average Precision \\
\hline
RQ4 & How does local claim-submission latency and throughput change as concurrent
submissions increase? &
Exploratory local Hardhat benchmark at five concurrency levels \\
\hline
RQ5 & Does auditor reputation align with historical agreement with finalized
consensus? &
Pearson-correlation study, executable after finalized voting history exists \\
\hline
\end{tabular}
}
\end{table}

\section{System Architecture and Trust Boundaries}
\label{sec:system_architecture}

\subsection{Hybrid Architecture}

Block-Insure uses a hybrid architecture because no single storage or execution
layer satisfies all requirements. Blockchain execution is appropriate for
shared lifecycle state, access control, duplicate markers, quorum decisions, and
financial settlement. It is unsuitable for storing large medical documents or
performing flexible statistical analysis. Conversely, an ordinary database is
efficient and searchable, but a database administrator can modify its contents.
The resulting design places compact, decision-critical commitments on-chain and
keeps detailed or high-volume information off-chain.

\begin{table}[htbp]
\centering
\caption{Architectural components and authority boundaries}
\label{tab:architecture_authority}
\resizebox{\textwidth}{!}{
\begin{tabular}{|p{3.0cm}|p{5.5cm}|p{4.5cm}|p{3.2cm}|}
\hline
\textbf{Component} & \textbf{Main Responsibility} &
\textbf{Authoritative For} & \textbf{Not Authoritative For} \\
\hline
React frontend & Role-specific interfaces and wallet interaction &
No persistent system state & Claim truth, roles, or settlement \\
\hline
Express backend and MongoDB & Authentication, profiles, metadata, registry
simulation, analytics, notifications, appeals, and orchestration &
Application metadata and research artifacts & Final on-chain claim state \\
\hline
IPFS through Pinata & Content-addressed storage of uploaded evidence bytes &
Retrieval by CID & Confidentiality or access control \\
\hline
\texttt{InsuranceManager} contract & Policies, claims, roles, duplicate markers,
oracle quorum, votes, reputation, registry root, and settlement &
Lifecycle and financial state & Truthfulness of off-chain source data \\
\hline
Oracle workers & Registry query, risk assessment, response hashing, and
confirmation submission & Their signed on-chain confirmations &
Independent truth when infrastructure is shared \\
\hline
\end{tabular}
}
\end{table}

\subsection{Technology Stack}

\begin{table}[htbp]
\centering
\caption{Principal implementation technologies}
\label{tab:technology_stack_methodology}
\begin{tabular}{|p{4.0cm}|p{4.8cm}|p{5.2cm}|}
\hline
\textbf{Layer} & \textbf{Technology} & \textbf{Purpose} \\
\hline
Blockchain & Solidity 0.8.28, Hardhat, OpenZeppelin & Contract implementation,
local execution, access control, pause, and reentrancy protection \\
\hline
Application API & Node.js, Express, ethers.js & Authentication, contract
interaction, orchestration, and analytics \\
\hline
Database & MongoDB with Mongoose & Profiles, evidence metadata, registries,
oracle logs, notifications, and research records \\
\hline
Evidence storage & Pinata HTTP API and IPFS & Content-addressed document storage \\
\hline
Oracle layer & Two Node.js workers using ethers.js and Axios & Polling
\texttt{OracleRequested} events and submitting confirmations \\
\hline
Presentation & React and Vite & User, administrator, auditor, and thesis-results
dashboards \\
\hline
\end{tabular}
\end{table}

\subsection{Data Placement}

\begin{table}[htbp]
\centering
\caption{Data placement and integrity rationale}
\label{tab:data_placement}
\resizebox{\textwidth}{!}{
\begin{tabular}{|p{4.1cm}|p{3.7cm}|p{7.4cm}|}
\hline
\textbf{Data} & \textbf{Location} & \textbf{Rationale and Limitation} \\
\hline
Roles, configuration, policies, core claims, status, and trust score &
Smart contract & Shared and auditable state; publicly visible and expensive to
store \\
\hline
Initial evidence hash and CID & Smart contract & Immutable commitment to the
submitted evidence reference; does not hide the CID \\
\hline
Raw evidence bytes & IPFS/Pinata & Avoids on-chain file storage; evidence should
be encrypted in a production medical system \\
\hline
Evidence metadata and evidence-chain links & MongoDB & Flexible append-style
tracking; detects inconsistency but remains alterable by a database
administrator \\
\hline
Registry records and detailed oracle responses & MongoDB & Searchable external
data simulation and detailed logs; correctness depends on the operator \\
\hline
Registry Merkle root and oracle response hash & Smart contract & Compact
integrity commitments; current prototype treats mismatches as audit signals,
not automatic rejection gates \\
\hline
Auditor votes and reputation & Smart contract & Transparent vote records and
reputation values; weighted consensus is calculated by the backend \\
\hline
Settlement parameters, transfer, and record & Smart contract & Enforced
financial calculation and payout \\
\hline
\end{tabular}
}
\end{table}

\section{Actors, Identity, and Authorization}
\label{sec:actors_identity}

\subsection{Actor Model}

The prototype supports policyholders, claim officers, administrators, oracle
workers, and auditors. A policyholder purchases policies, submits claims, and
may appeal a rejection. A claim officer requests oracle verification and makes
eligible claim decisions. An administrator manages configuration, roles,
registry commitments, appeals, settlements, and closure. Oracle workers submit
external-verification confirmations. Auditors inspect histories and vote on
claims in manual-review or oracle-failed states.
Weighted \texttt{NEEDS\_MORE} consensus is treated as non-final: reputation
updates and claim-decision gating remain open until additional evidence and
votes produce a valid-or-invalid consensus.

\begin{table}[htbp]
\centering
\caption{On-chain roles and security boundaries}
\label{tab:onchain_roles}
\resizebox{\textwidth}{!}{
\begin{tabular}{|p{3.3cm}|p{8.5cm}|p{4.2cm}|}
\hline
\textbf{Role} & \textbf{Authorized Operations} & \textbf{Important Boundary} \\
\hline
\texttt{ADMIN\_ROLE} & Configuration, role and package management, registry-root
updates, timeout resolution, appeal finalization, reputation updates,
settlement, closure, funding, and withdrawal & Highly privileged; final admin
cannot be removed \\
\hline
\texttt{CLAIM\_OFFICER\_ROLE} & Request oracle verification, send to manual
review, approve, and reject & Cannot settle, withdraw, configure, or manage
roles \\
\hline
\texttt{ORACLE\_ROLE} & Submit one confirmation per oracle request & Confirmation
does not prove that its external source is truthful \\
\hline
\texttt{AUDITOR\_ROLE} & Cast one vote on an eligible claim & Cannot finalize
consensus or directly update reputation \\
\hline
\texttt{EMERGENCY\_ROLE} & Pause the contract & Cannot unpause; only an
administrator can restore operation \\
\hline
\end{tabular}
}
\end{table}

The backend separately stores the application roles
\texttt{USER}, \texttt{ADMIN}, \texttt{AUDITOR}, and \texttt{ORACLE}. Backend
middleware protects API routes, while Solidity roles protect transactions.
These two role systems must be configured consistently. A mismatch can expose a
page without granting the required contract permission, or grant a contract
permission without exposing the corresponding page.

\subsection{Wallet-Signature Authentication}

Authentication is passwordless. The backend creates a random 16-byte nonce and
returns the message
\texttt{Block-Insure login nonce: <nonce>}. The wallet signs this message, the
backend recovers the signer with \texttt{verifyMessage}, and login succeeds only
when the recovered address equals the claimed wallet. The nonce is cleared after
successful use. A seven-day JWT containing a unique \texttt{jti} is then issued.
During logout, the \texttt{jti} is placed in a MongoDB revocation collection
until token expiry; authentication middleware checks both revocation status and
the user's current database role.

\section{On-Chain Claim-Processing Method}
\label{sec:onchain_method}

\subsection{Claim Lifecycle}

The \texttt{InsuranceManager} contract represents each claim using the states in
Table~\ref{tab:claim_states}. State-changing functions validate the current
state and caller role before performing a transition.

\begin{table}[htbp]
\centering
\caption{Claim lifecycle states and principal transitions}
\label{tab:claim_states}
\resizebox{\textwidth}{!}{
\begin{tabular}{|c|p{3.2cm}|p{10.2cm}|}
\hline
\textbf{Code} & \textbf{State} & \textbf{Meaning and Principal Transition} \\
\hline
0 & \texttt{SUBMITTED} & Temporary creation state inside
\texttt{submitClaim} \\
\hline
1 & \texttt{DUPLICATE\_CHECKED} & Passed deterministic duplicate checks; may
enter \texttt{ORACLE\_PENDING} \\
\hline
2 & \texttt{FRAUD\_FLAGGED} & Duplicate marker detected; may enter
\texttt{MANUAL\_REVIEW} \\
\hline
3 & \texttt{ORACLE\_PENDING} & Awaiting confirmations; becomes verified,
failed, or timed out \\
\hline
4 & \texttt{ORACLE\_VERIFIED} & Verified by quorum; may be approved or rejected \\
\hline
5 & \texttt{ORACLE\_FAILED} & Failed quorum or timeout; may enter manual review \\
\hline
6 & \texttt{MANUAL\_REVIEW} & Human-review state; may be approved or rejected \\
\hline
7 & \texttt{APPROVED} & Eligible for administrator-only settlement \\
\hline
8 & \texttt{REJECTED} & May be appealed, reopened, finalized, or closed after
the appeal window \\
\hline
9 & \texttt{SETTLED} & On-chain payout completed; may be closed \\
\hline
10 & \texttt{CLOSED} & Terminal administrative closure state \\
\hline
\end{tabular}
}
\end{table}

The ordinary successful path is
\begin{equation}
\texttt{DUPLICATE\_CHECKED}
\rightarrow \texttt{ORACLE\_PENDING}
\rightarrow \texttt{ORACLE\_VERIFIED}
\rightarrow \texttt{APPROVED}
\rightarrow \texttt{SETTLED}
\rightarrow \texttt{CLOSED}.
\end{equation}
The contract also supports fraud escalation, oracle failure, rejection, a
single claimant appeal, reopening for a fresh oracle request, and explicit
closure. A claim officer can reject any existing claim that is neither settled
nor closed; this broad authority is a deliberate prototype simplification and
should be narrowed in a production workflow.

\subsection{Submission Validation and Duplicate Detection}

A submission is accepted only when the policy exists, belongs to the caller, is
active and unexpired, the incident date lies inside the policy period and is not
in the future, the amount is positive and within coverage, required text and
hash fields are present, and the configurable per-policy claim limit has not
been reached. The default maximum is five claims per policy.

After validation, the contract tests three global or claimant-specific
duplicate markers:
\begin{equation}
D =
D_{\mathrm{document}}
\lor D_{\mathrm{invoice}}
\lor D_{\mathrm{wallet,date,type}}.
\label{eq:duplicate_rule}
\end{equation}
If \(D\) is true, the claim becomes \texttt{FRAUD\_FLAGGED}. Otherwise, the
document hash, invoice hash, and wallet/date/type fingerprint are registered as
used. Fraud-flagged submissions return before registering their new markers, so
an invalid submission does not poison unused identifiers for later legitimate
claims. These checks detect exact reuse; they do not detect visually similar
documents, collusion, or a newly hashed fabricated document.

\subsection{Deterministic On-Chain Trust Score}

For a clean submission, the contract computes a deterministic trust score.
Higher values indicate greater trust, which is the opposite direction from the
off-chain Bayesian fraud probability.

\begin{equation}
\begin{split}
T_{\mathrm{raw}}={}&15I_{\mathrm{activePolicy}}
+15I_{\mathrm{incidentInPeriod}}
+15I_{\mathrm{coveredAmount}}
+15I_{\mathrm{documentPresent}}\\
&+10I_{\mathrm{newDocument}}
+10I_{\mathrm{newInvoice}}
+10I_{\mathrm{newWalletDateType}}
+25I_{\mathrm{oracleVerified}},
\end{split}
\label{eq:onchain_trust_score}
\end{equation}
\begin{equation}
T = \min(T_{\mathrm{raw}},100).
\end{equation}

The pre-oracle maximum is 90, and successful oracle verification increases it
to the capped value of 100. Scores \(T\geq80\), \(50\leq T<80\), and \(T<50\)
map to \texttt{LOW}, \texttt{MEDIUM}, and \texttt{HIGH} risk labels,
respectively. \texttt{FRAUD\_FLAGGED} and \texttt{ORACLE\_FAILED} override the
numeric label. Because submission validation and duplicate gating already
require all seven pre-oracle conditions for a clean claim, every clean claim
begins at 90 and a verified claim reaches 100. The on-chain score is therefore
best interpreted as a deterministic lifecycle-confidence indicator rather than
a calibrated or highly granular fraud-risk model.

\subsection{Oracle Request, Quorum, and Timeout}

A claim officer or administrator creates an oracle request only for a
\texttt{DUPLICATE\_CHECKED} claim. The query commitment is
\begin{equation}
Q_h=\mathrm{Keccak256}\left(
\mathrm{abi.encodePacked}(claimId,policyId,wallet,hospitalId,invoiceHash,documentHash)
\right).
\label{eq:oracle_query_hash}
\end{equation}

Each address with \texttt{ORACLE\_ROLE} may confirm once. After the configurable
threshold is reached, the request succeeds only when verified confirmations
strictly outnumber failed confirmations:
\begin{equation}
Verified_{\mathrm{final}} =
\left(N_{\mathrm{verified}} > N_{\mathrm{failed}}\right).
\label{eq:oracle_majority}
\end{equation}

\begin{table}[htbp]
\centering
\caption{Default two-oracle quorum truth table}
\label{tab:oracle_truth_table}
\begin{tabular}{|c|c|c|c|}
\hline
\textbf{Oracle 1} & \textbf{Oracle 2} & \textbf{Final Result} &
\textbf{Claim State} \\
\hline
Verified & Verified & Verified & \texttt{ORACLE\_VERIFIED} \\
\hline
Verified & Failed & Failed & \texttt{ORACLE\_FAILED} \\
\hline
Failed & Verified & Failed & \texttt{ORACLE\_FAILED} \\
\hline
Failed & Failed & Failed & \texttt{ORACLE\_FAILED} \\
\hline
\end{tabular}
\end{table}

The default threshold is two and the default timeout is 50 blocks. After the
timeout, an administrator can finalize the pending request as failed. The
threshold can be configured above the number of available oracle workers;
therefore, incorrect configuration may intentionally or accidentally force
requests to time out.

The contract stores one \texttt{resultHash}, \texttt{riskLevel}, and
\texttt{remarks} field in the request. These fields are overwritten by each
confirmation, so the final stored response hash belongs to the confirmation
that reaches quorum; it is not an aggregate commitment to all oracle responses.
Per-oracle detailed responses remain in MongoDB.

\subsection{Auditor Voting and Reputation}

Auditors can vote \texttt{VALID}, \texttt{INVALID}, or \texttt{NEEDS\_MORE} on
claims in \texttt{MANUAL\_REVIEW} or \texttt{ORACLE\_FAILED}. The vote itself
is stored on-chain. The weighted consensus is calculated by the backend:
\begin{equation}
w_i=\frac{\mathrm{clamp}(R_i,0,100)}{100},
\qquad
W_c=\sum_{i\in V_c}w_i,
\label{eq:weighted_consensus}
\end{equation}
where \(R_i\) is auditor reputation and \(V_c\) is the set of votes for category
\(c\). A unique maximum becomes the consensus; a tie is not finalized.
Auditors agreeing with consensus receive \(+2\) reputation and others receive
\(-1\), clamped to \([0,100]\).

This consensus and reputation-update procedure is \emph{not} autonomously
executed by the smart contract. An authenticated administrator triggers backend
finalization, and the backend submits administrator-signed reputation-update
transactions. Consequently, on-chain votes are transparent, but the current
finalization process still trusts the backend and administrator. The contract
also permits an administrator to set any auditor reputation from 0 to 100
directly, independently of the backend consensus calculation.

\subsection{Settlement and Reserve Intelligence}

Settlement uses a configurable deductible and insurer share. With claim amount
\(C\), deductible rate \(r_d\), deductible cap \(d_{\max}\), and insurer share
\(r_i\):
\begin{equation}
d=\min(Cr_d,d_{\max},C),
\qquad
P_{\mathrm{insurer}}=(C-d)r_i,
\qquad
P_{\mathrm{claimant}}=C-P_{\mathrm{insurer}}.
\label{eq:settlement_formula}
\end{equation}
The defaults are \(r_d=10\%\), \(d_{\max}=0.02\) ETH, and \(r_i=80\%\).

\begin{table}[htbp]
\centering
\caption{Worked settlement values using default parameters}
\label{tab:settlement_examples}
\begin{tabular}{|r|r|r|r|}
\hline
\textbf{Claim (ETH)} & \textbf{Deductible (ETH)} &
\textbf{Insurer Pays (ETH)} & \textbf{Claimant Responsibility (ETH)} \\
\hline
0.01 & 0.001 & 0.0072 & 0.0028 \\
\hline
0.05 & 0.005 & 0.0360 & 0.0140 \\
\hline
0.10 & 0.010 & 0.0720 & 0.0280 \\
\hline
0.50 & 0.020 & 0.3840 & 0.1160 \\
\hline
\end{tabular}
\end{table}

Only an administrator can settle an \texttt{APPROVED} claim. The contract
checks available balance, records the settlement, changes state before the
external transfer, uses \texttt{nonReentrant}, and emits a low-reserve warning
when the remaining balance is below the default 0.1 ETH threshold. The backend
also calculates open exposure and solvency labels for the dashboard. These
reserve calculations are advisory: \texttt{withdrawExcess} does not enforce
outstanding-liability coverage, so a privileged administrator can withdraw
funds that may be needed for future claims.

\section{Evidence and Registry Integrity}
\label{sec:evidence_registry}

\subsection{Evidence Upload and Verification}

The backend accepts a document in memory, calculates
\(\mathrm{SHA256}(\mathrm{fileBytes})\), uploads the bytes to Pinata/IPFS, and
stores the SHA-256 value, CID, MIME type, uploader, and document type in
MongoDB. The initial claim stores the document hash and CID on-chain. An auditor
can later hash a supplied file and compare it with the on-chain value.

IPFS content addressing detects byte changes because modified content receives a
different CID. It does not provide confidentiality. Unless evidence is
encrypted before upload or protected by a private retrieval layer, anyone who
obtains the CID may be able to retrieve the document.

\subsection{Off-Chain Evidence Chain}

Additional evidence metadata is linked in MongoDB. For evidence item \(E_i\),
the implementation calculates
\begin{equation}
H_i=\mathrm{SHA256}\left(
\mathrm{JSON}(claimId,i,H_{i-1},sha256Hash,CID,documentType,uploader)
\right),
\label{eq:evidence_chain}
\end{equation}
where \(H_{-1}=\texttt{GENESIS}\). Verification recomputes each link and checks
the index and previous hash. This detects inconsistent edits under normal
operation. Because both the records and chain hashes are stored in the same
database, a database administrator could rewrite the chain consistently.
Anchoring each head hash on-chain would be required for stronger tamper
resistance.

\subsection{Synthetic Registry and Merkle Commitment}

The synthetic primary registry contains 72 healthcare records and fields such as
hospital identity, license status, patient hash, treatment type, dates, bill
range, invoice identity, statuses, fraud scenario, and previous-claim count. A
second registry collection is seeded separately for Oracle 2 and intentionally
marks one otherwise legitimate invoice as used.

Records are sorted by normalized invoice hash and invoice number. Each canonical
record is serialized in a fixed field order and hashed:
\begin{equation}
L_i=\mathrm{SHA256}\left(\mathrm{UTF8}(\mathrm{JSON}(r_i))\right).
\label{eq:merkle_leaf}
\end{equation}
For child hashes \(a\) and \(b\), the actual implementation removes their
\texttt{0x} prefixes, concatenates the hexadecimal \emph{text}, and hashes its
UTF-8 representation:
\begin{equation}
P=\mathrm{SHA256}\left(
\mathrm{UTF8}(\mathrm{hex}(a)\parallel\mathrm{hex}(b))
\right).
\label{eq:merkle_parent}
\end{equation}
When a level has an odd number of nodes, its final node is paired with itself.
The administrator commits the resulting root, snapshot timestamp, and block
number to the contract.

An inclusion proof can demonstrate that a record belongs to the committed
snapshot. It cannot demonstrate that the record is truthful. Furthermore, the
current oracle records whether its local root matches the on-chain root, but
root mismatch does not automatically change the oracle's \texttt{verified}
decision. Thus, Merkle commitment is an implemented audit mechanism rather than
an enforced oracle-security gate in this prototype.

\section{Oracle Verification and Fraud-Risk Method}
\label{sec:oracle_fraud_method}

\subsection{Oracle Worker Procedure}

Each oracle worker polls \texttt{OracleRequested} events, reads the claim,
queries its configured registry snapshot, obtains a rule-based comparison,
Bayesian assessment, and Merkle proof, then constructs a detailed response. The
response commitment is
\begin{equation}
R_h=\mathrm{Keccak256}\left(
\mathrm{UTF8}(\mathrm{JSON.stringify}(response))
\right).
\label{eq:oracle_response_hash}
\end{equation}
The worker submits the Boolean verification result, \(R_h\), risk level, and
remarks on-chain, waits for the receipt, and attempts to store the full response
and elapsed time in MongoDB. Oracle 2 adds a configurable default 3000 ms
simulated network-propagation delay.

Both workers use different registry collections and different oracle wallets,
but they still share one backend and MongoDB deployment. Therefore, they
simulate different oracle views without providing production-grade
independence.

\subsection{Rule-Based Registry Decision}

The registry verifier compares hospital ID, invoice hash, claim amount,
treatment type, record status, invoice status, hospital-license status, and the
synthetic fraud label. Date consistency is a non-blocking warning using a
\(\pm30\)-day tolerance. Expected treatment bill range also allows a 15\%
tolerance. The worker verifies a record only when no blocking comparison fails
and the Bayesian recommendation is not
\texttt{REJECT\_ORACLE\_VERIFICATION}.

The synthetic \texttt{fraudLabel} is directly checked by this rule-based
verifier, but it is not used as an input factor by the Bayesian model. This
distinction is important: a real external registry would not ordinarily expose
a reliable ground-truth fraud label, so production verification must replace
this simulation field with independently verifiable evidence.

\subsection{Bayesian Fraud Probability}

The off-chain scorer applies Naive Bayes to the set \(A\) of active evidence
factors. Examples include clean registry match, invalid status, used invoice,
cancelled record, blacklisted license, bill anomaly, and repeated-claim pattern.
Only active factors are multiplied; inactive-factor complements are not included
in the current implementation.

\begin{equation}
\ell_F=\log P(F)+\sum_{j\in A}\log P(x_j=1\mid F),
\qquad
\ell_L=\log P(L)+\sum_{j\in A}\log P(x_j=1\mid L),
\label{eq:bayes_logs}
\end{equation}
\begin{equation}
P(F\mid A)=
\frac{\exp(\ell_F-m)}
{\exp(\ell_F-m)+\exp(\ell_L-m)},
\qquad
m=\max(\ell_F,\ell_L).
\label{eq:bayes_posterior}
\end{equation}

Likelihoods are trained from synthetic registry records with Laplace smoothing:
\begin{equation}
P(x_j=1\mid C)=
\frac{N(x_j=1,C)+1}{N(C)+2}.
\label{eq:laplace}
\end{equation}
The fraud label is used as the training target, not as a Bayesian evidence
factor. Some runtime mismatch factors have no positive examples in the
synthetic training data and therefore receive only smoothed likelihoods; their
probabilities should not be interpreted as learned real-world effects.

Numeric anomaly signals use sample variance:
\begin{equation}
s^2=\frac{\sum_{i=1}^{N}(x_i-\bar{x})^2}{N-1},
\qquad
z=\frac{x-\bar{x}}{s},
\label{eq:zscore}
\end{equation}
with an anomaly threshold of \(|z|\geq2\). Fraud probability below 35 is
\texttt{LOW}, 35--69 is \texttt{MEDIUM}, and at least 70 is \texttt{HIGH}.
A score of at least 85 or any blocking comparison produces the rejection
recommendation.

\begin{table}[htbp]
\centering
\caption{Distinction between the two implemented risk measures}
\label{tab:risk_measure_comparison}
\begin{tabular}{|p{3.4cm}|p{5.4cm}|p{5.4cm}|}
\hline
\textbf{Property} & \textbf{On-Chain Trust Score} &
\textbf{Off-Chain Bayesian Probability} \\
\hline
Direction & Higher is safer & Higher is more suspicious \\
\hline
Inputs & Policy validity, amount, evidence presence, duplicate markers, oracle
result & Active registry, status, anomaly, and repeat-pattern factors \\
\hline
Computation & Deterministic Solidity rules & Trained likelihoods and log-space
posterior \\
\hline
Primary role & Auditable lifecycle signal & Explainable fraud-risk
recommendation \\
\hline
\end{tabular}
\end{table}

\section{Security Controls, Assumptions, and Limitations}
\label{sec:method_security_limitations}

\begin{table}[htbp]
\centering
\caption{Implemented controls and residual trust assumptions}
\label{tab:controls_assumptions}
\resizebox{\textwidth}{!}{
\begin{tabular}{|p{4.0cm}|p{5.7cm}|p{5.8cm}|}
\hline
\textbf{Control} & \textbf{Implemented Protection} &
\textbf{Residual Limitation} \\
\hline
Role-based access control & Restricts privileged contract and API operations &
Backend and on-chain roles require synchronization \\
\hline
Final-admin protection and emergency pause & Prevents removal of the final
administrator; emergency role can pause but not unpause & Administrator remains
a powerful trusted party \\
\hline
Reentrancy guard and checks-effects-interactions & Protects settlement and
withdrawal transfers & Does not prevent privileged misuse of valid functions \\
\hline
JWT revocation & Rejects logged-out token identifiers until expiry & Stolen
wallet keys still authenticate the attacker \\
\hline
IP and wallet submission limits & Defaults to 5 requests per 15 minutes and 3
wallet attempts per 24 hours & Direct contract calls bypass backend limits;
on-chain per-policy cap is the backstop \\
\hline
Encrypted evidence, document hash, CID, evidence chain, and Merkle root &
Supports confidentiality, integrity checking, and auditability & Browser-held
decryption keys can be lost; off-chain chain and registry remain operator-managed \\
\hline
Two-oracle quorum and timeout & Requires two confirmations, locally verifies
Merkle proofs, enforces the committed root, and fails split decisions
conservatively & Workers still share prototype infrastructure \\
\hline
Reserve intelligence and warning & Exposes liabilities, emits low-reserve
warning, and prevents withdrawals from consuming approved liabilities &
Prototype funding remains manually administered \\
\hline
\end{tabular}
}
\end{table}

The backend signs privileged transactions with one environment-configured
\texttt{ADMIN\_PRIVATE\_KEY}. A production deployment should use multisignature
custody, timelocks, separation of treasury and operational signers, monitoring,
and key rotation. The optimized monolithic contract has a deployed runtime size
of 24,340 bytes, which is 236 bytes below the 24,576-byte EIP-170 limit.
Hardhat's unlimited-size bypass is disabled and a regression test protects the
boundary.

\section{Experimental Design and Reproducibility}
\label{sec:experimental_design}

\begin{table}[htbp]
\centering
\caption{Experimental protocols used to evaluate the artifact}
\label{tab:experiment_protocols}
\resizebox{\textwidth}{!}{
\begin{tabular}{|p{3.2cm}|p{5.4cm}|p{4.5cm}|p{3.8cm}|}
\hline
\textbf{Study} & \textbf{Independent Variable or Input} &
\textbf{Measured Output} & \textbf{Control or Reproducibility Mechanism} \\
\hline
Contract verification & Functions, roles, valid and invalid transitions &
Pass/fail count and code coverage & Hardhat fixtures and deterministic local
chain \\
\hline
Gas comparison & Registry size \(N\in\{10,25,50,100,200,500\}\) and storage
method & Receipt gas used and percentage saving & Fresh \texttt{GasTestContract}
for each measurement \\
\hline
Fraud-model evaluation & Deterministically seeded 72-record registry &
Threshold, confusion matrix, baselines, AUC, AP, and risk distribution &
Stable sort; four-of-five training split; every fifth record held out \\
\hline
Concurrency benchmark & Concurrent submissions
\(C\in\{1,5,10,20,50\}\) & Successes, failures, average and p95 latency, and
wall-clock throughput & One policy per claim; policy setup excluded from timing \\
\hline
Auditor study & Finalized voting history & Reputation/accuracy Pearson
correlation & Refuses to generate a result without finalized sessions \\
\hline
\end{tabular}
}
\end{table}

The generated JSON, CSV, and PNG artifacts are retained under
\texttt{backEnd/evaluation-results}, while gas measurements are retained in
\texttt{contracts/gas-comparison-results.csv}. The next chapter reports the
measured results and discusses their validity.

\section{Summary}

Block-Insure was implemented as a hybrid research artifact in which the smart
contract enforces shared claim and settlement state while off-chain services
provide storage, external-data simulation, analytics, and user interaction. The
methodology combines deterministic on-chain controls with explainable off-chain
risk scoring and explicitly identifies where trust remains centralized. This
separation allows the experimental results to be interpreted as evidence of
prototype correctness and feasibility, rather than as proof of production
decentralization or real-world fraud-detection accuracy.
